\section[Appendix B {}- Running the Test Cases]{Appendix B - Running the Test Cases}
{
This section describes how to run the Dyninst test cases. The primary purpose of the test cases is to verify that the
API has been installed correctly (and for use in regression testing by the developers of the Dyninst library). The
code may also be of use to others since it provides a fairly complete example of how to call most of the API methods.
The test suite consists of mutator programs and their associated mutatee programs.

{
To compile the test suite, type make in the appropriate platform specific directory under dyninst/testsuite. To run,
execute runTests. Each test will be executed and the result (PASSED/FAILED/CRASHED) printed.}

{
Test mutators are run by the test_driver executable (test_driver.exe on Windows). The test_driver loads a mutator
test from a shared object and runs it on a test mutatee. A single run of the test_driver may execute multiple tests
(depending on parameters passed), and each test may execute multiple times with different parameters and on different
mutatees.

{
Dyninst's test space can be very large. Each mutatee can be run under different tests, compiled
by different compilers, and run with different parameters. For example, one point in this space would be the test1
mutatee being run under under test1_13, when compiled with the g++ compiler, and in attach mode. When run without
any options, the test_driver will run all test combinations that are valid on the current platform. Many of the
options that are passed to test_driver can be used to limit the test space that it runs in.}

{
In order to prevent a crashing test from stopping the test_driver from running subsequent tests, test_driver can be
run under a wrapper application, runTests. The runTests wrapper invokes the test_driver with the any arguments that
were passed to runTests. It will watch the test_driver process, and if test_driver exits with a fault it will print
an appropriate error message and restart the test_driver on the next test.}

{
It is generally recommended that runTests be used when running a large sequence of tests, and test_driver be used when
debugging issues with a single test.}

{
The test_driver and runTests applications can be invoked with the following list of arguments. Most arguments are
used to limit the space of tests that the testsuite will run. For example, to run the above test1_13 example, you
could use the following command line:}

{
test_driver –run test1_13 –mutatee test1.mutatee_g++ -attach}


\bigskip

{
{}-attach

{
Only run tests that attach to the mutatees.}

{
{}-create}

{
Only run tests that create mutatees.}

{
{}-rewriter}

{
Only run tests that rewrite mutatees.}

{
{}-staticlink}

{
Run rewriter tests that use statically linked mutatees.}

{
{}-dynamiclink}

{
Run rewriter tests that use dynamically linked mutatees.}

{
{}-allmode}

{
Run tests for all modes (create, attach, rewriter on statically linked binaries, rewriter on dynamically linked
binaries).}

{
{}-gcc, -g++, -pgcc, -pgCC, -icc, -icpc}

{
Run tests on mutatees built with the specified compiler.}

{
{}-noclean}

{
Don't remove rewritten mutatees after running rewriter tests.}

{
{}-all}

{
Run tests for all possible combinations of unoptimized mutatees.}

{
{}-none, -low, -high, -max}

{
Only run tests for mutatees of the given optimization level.}

{
{}-allopt}

{
Run tests for all mutatee optimization levels.}

{
{}-full}

{
Run tests for all possible combinations of mutatees, including all optimization levels. Requires make all to build the
optimized mutatees.}

{
{}-dyninst, -symtab, -instruction, -proccontrol, -stackwalker}

{
Only run tests for the specified component.}

{
{}-allcomp}

{
Run tests for all components.}

{
{}-32, -64}

{
Only run tests for 32-bit or 64-bit mutatees. This option is only valid on platforms such as AMD64/Linux and
PowerPC/Linux where both 32-bit and 64-bit build environments are available.}

{
{}-pic}

{
Only run tests for mutatees compiled as position-independent code. Default is non-PIC mutatees. –all and –full
include both.}

{
{}-sp, -mp, -st, -mt}

{
ProcControl-specific: run tests in single process, multiprocess, single thread, multithread modes, respectively.}

{
{}-j n}

{
This option spawns up to n test_driver instances from a given runTests invocation. This option is highly recommended
for large test runs.}

{
{}-hosts host1 [host2 ... hostn]}

{
In conjunction with the –j option above, will distribute tests over the hosts host1…hostn. The hosts must share a
filesystem with the machine from which runTests is being run, and ssh must be configured to allow password-less
authentication to those hosts. Note that –j controls the total number of test_driver instances, not the number per
host, so you will need to use a –j N at least equal to the number of hosts you wish to use.}

{
{}-mutatee <mutatee_name>}


\bigskip

{
Only run tests that use the specified mutatee name. Only certain mutatees can be run with certain tests. The
primary test number specifies which mutatees it can be run with. For example, all of the test1_* tests can be run
with the test1.mutatee_* mutatees, and all of the test2_* tests can be run with the test2.mutatee_* mutatees.}


\bigskip

{
{}-run <subtest> <subtest> …}

{
Only runs the specific sub-tests listed. For example, to run sub-test case 4 of test2 you would enter test_driver
–run test2_4.}

{
{}-test}

{
Alias for –run.}

{
{}-log}

{
Print more detailed output, including messages generated by the tests. Without this option the testsuite will capture
and hide any messages printed by the test, only showing a summary of whether the test passed or failed. By default,
output is sent to stdout.}

{
{}-logfile <filename>}

{
Send output from the –log option to the given filename rather than to stdout.

{
{}-verbose}

{
Enables test suite debugging output. This is useful when trying to track down issues in the test suite or tests.}
